\section{Sistema de Ecuaciones Lineales.}
Objetivos:
\begin{itemize}
    \item Estudiar sistemas de ecuaciones lineales.
    \item Aprender a utilizar el método de la reducción gaussiana para resolver sistemas de ecuaciones algebraicas.
    \item Aprender a utilizar el método de la reducción de Gauss-Jordan para resolver sistemas de ecuaciones algebraicas.
    \item Aprender los sistemas homogéneos y las características de su solución.
    \item Relacionar matrices y vectores con sistemas de ecuaciones.
    \item Determinar condiciones para la existencia de soluciones de sistemas de ecuaciones homogéneos y no homogéneos.
\end{itemize}
\subsection{Introducción a los sistemas lineales.}

\begin{defin}
Se define una ecuación lineal de $n$ variables sobre un cuerpo $(K, +, ·)$
como una igualdad de la forma: 
\[a_1\cdot x_1 + a_2\cdot x_2 + \dots  + a_n\cdot x_n = b,\]
donde cada $x_j$ es variable incógnita de la ecuación, cada $a_j$ es coeficiente numérico (en $K$)
para la $j-$ésima variable $(x_j)$ y $b$ es la constante (en $K$) que resulta de la suma ponderada (por $a_j$) de las variables $x_j$. De manera abreviada
\[\sum_{j=1}^n a_j \cdot x_j=b\]
\end{defin}
Su resolución consiste en determinar todas y cada una de las $n-$uplas $(s_1, s_2,\dots, s_n )$ de valores
$s_j$ asignados a cada variable $x_j$ para satisfacer la igualdad, es decir 
\[\sum_{j=1}^n a_j \cdot s_j=b\]
Luego, el conjunto solución se anotaría:

\[S_E=\left\{(s_1, s_2,\dots, s_n )\in K^n \colon \sum_{j=1}^n a_j \cdot s_j=b\right\}\]

\begin{ejem}
$3x_1 - 2x_2 + 7x3 - x_4 = 9$ es una ecuación lineal de 4 variables.
$(-1, 2, 3, 5 )$ es \textbf{una} solución o un elemento del conjunto solución puesto que : $3(-1) - 2(2) + 7(3) -(5)$ resulta igual a 9.

Verifique usted que $(2, 3, -2, -23)$ es otra solución para nuestra ecuación.
Note que hay una infinidad de cuádruplas en la solución $S_E$ de nuestra ecuación.
\end{ejem}

Es muy frecuente que nos interese resolver problemas que exigen satisfacer varias
condiciones (ecuaciones) simultáneas por parte de un cierto número $n$ de variables
(incógnitas). En este caso, se habla de resolver una familia de $m$ ecuaciones para esas
incógnitas. Como la resolución debe ser simultánea, el conjunto solución (de $n-$uplas) para
el sistema de ecuaciones resultará coincidente con la intersección de los $m$ conjuntos
solución para cada una de esas $m$ ecuaciones.

\begin{defin}
Un \textbf{sistema de $m$ ecuaciones con $n$ incógnitas} es un sistema que se escribe de la siguiente forma:
\[\begin{array}{ccccccc}
        a_{11}x_1&+ a_{12} x_2&+ \dots &+ a_{1j}x_j &+ \dots &+ a_{1n}x_n=& b_1\\
     a_{21}x_1&+ a_{22} x_2&+ \dots &+ a_{2j} x_j&+\dots &+  a_{2n}x_n=& b_2\\
    \vdots & \vdots &  & \vdots  &  &\dots  & \vdots\\
     a_{m1}x_1&+ a_{m2} x_2&+ \dots&+  a_{mj}x_j&+ \dots&+  a_{mn}x_n=& b_m
\end{array}\]
\end{defin}
\begin{rmk}
Note que, respecto de los coeficientes numéricos $a_{ij}$, $i$ nos indica que es coeficiente de la $i-$ésima ecuación y que $j$ nos indica que este coeficiente pondera a la $j-$ésima incógnita.
\end{rmk}
%\begin{ejem}

%Un \textbf{sistema de dos ecuaciones} y dos incógnitas $x$ e $y$ es un sistema de la forma:
%\begin{align}\label{sistema}
% a_{11}x + a_{12}y &= b_1,\\
%a_{21}x + a_{22}y &= b_2,   \nonumber
%\end{align}
%donde $ a_{11}, a_{12},a_{21},a_{22}, b_1$ y $
%b_2$ son n\'umeros dados. 

%Una \textbf{solución al sistema} \ref{sistema} es un par de números, denotados por $(x,y)$, que satisface las igualdades en \ref{sistema}. 

%\end{ejem}

Si interpretamos cada una de las ecuaciones del sistema como una l\'inea recta, la soluci\'on de dicho sistema puede pensarse como la intersecci\'on de las dos rectas en un espacio $n-$dimensional.

Un sistema de ecuaciones lineales tiene ya sea:
\begin{enumerate}
    \item Una solución (el sistema es consistente).
    \item Infinitas soluciones (el sistema es consistente).
    \item Ninguna solución (el sistema es inconsistente).
\end{enumerate}
\begin{ejem}
\begin{enumerate}
    \item Resolver el sistema 
\begin{align*}
    3x -2y &= 4,\\
    5x + 2y &= 12.
\end{align*}
Sumando ambas ecuaciones tenemos que $8x=16$, luego $x=2$. Reemplazando en la primera ecuaci\'on $3\cdot 2-2y=4$, por lo tanto $2=2y$, obteniendo $y=1$. Finalmente la solución única de este sistema de ecuaciones es el vector $(2, 1)$.

Geom\'etricamente las ecuaciones representan dos rectas que se
intersectan en un punto $(2, 1)$.
    
    \item Resolver el sistema 
    \begin{align*}
    x -y &= 7,\\
    2x -2y &= 14.
\end{align*}
Notemos que la segunda ecuaci\'on puede obtenerse multiplicando la primera por dos, por lo tanto las ecuaciones representan la misma recta. Cada punto de la recta es una solución de este sistema. Por lo tanto, este sistema de ecuaciones tiene infinitas soluciones. Como $x -y = 7$, entonces $x = y -7$, luego las soluciones $(x,y)$ son de la forma $(y-7,y)$, obteniendo una soluci\'on por cada $y\in \R$. Por ejemplo, para $y=0$ tenemos el punto $(-7,0)$.

    \item Resolver el sistema 
    \begin{align*}
    x -y &= 7,\\
    2x -2y &= 13.
\end{align*}
Si asumimos que existe una soluci\'on, al restar ambas ecuaciones obtenemos $0=1$, lo que claramente es una contradicci\'on, por lo tanto, este sistema de ecuaciones no puede tener solución. En este caso, las ecuaciones representan dos rectas paralelas que no se intersectan.
\end{enumerate}
\end{ejem}

\begin{defin}
Si contamos con dos sistemas de ecuaciones que al menos coinciden en que sus incógnitas son las mismas, es decir, 
\[\sum_{j=1}^n a_j \cdot x_j=b_i,\]
con $1 \leq  i \leq  m$ y
\[\sum_{j=1}^n \alpha_j \cdot x_j=\beta_i\]
con $1 \leq  i \leq  p$ ; diremos que ellos son \textbf{equivalentes} cuando
y sólo cuando su conjunto solución sea el mismo.
\end{defin}





\subsection{Métodos directos para resolver sistemas lineales.}

En en caso particular de un sistema de $3\times 3$ podemos resolver dicho sistema siguiendo la siguiente estrategia:
\begin{enumerate}
    \item Dividir la primera ecuación, entre una constante, para hacer el coeficiente de $x_1$ igual a 1.
    \item Multiplicar la primera ecuación por las constantes adecuadas de modo que al sumarla a la segunda y tercera ecuaciones, respectivamente, se eliminen los términos en $x_1$.
    \item Dividir la segunda ecuación entre una constante, para hacer el coeficiente de $x_2$ igual a 1 y después usarla para eliminar los términos en $x_2$ de la primera y tercera ecuaciones.
    \item Dividir la tercera ecuación entre una constante, para hacer elcoeficiente de $x_3$ igual a 1 y después usarla para eliminar los términos de $x_3$ de la primera y segunda ecuaciones.
\end{enumerate}



\subsection{Sistemas de $m$ Ecuaciones con $n$ Incógnitas}
Al sistema 
\[\begin{array}{ccccccc}\label{nsistema}
        a_{11}x_1&+ a_{12} x_2&+ \dots &+ a_{1j}x_j &+ \dots &+ a_{1n}x_n=& b_1\\
     a_{21}x_1&+ a_{22} x_2&+ \dots &+ a_{2j} x_j&+\dots &+  a_{2n}x_n=& b_2\\
    \vdots & \vdots &  & \vdots  &  &\dots  & \vdots\\
     a_{m1}x_1&+ a_{m2} x_2&+ \dots&+  a_{mj}x_j&+ \dots&+  a_{mn}x_n=& b_m
\end{array}\]
$m$ ecuaciones con $n$ incógnitas se puede representar matricialmente de la forma:
\[\left(\begin{array}{cccc}
     a_{11}& a_{12} & \dots  & a_{1n}\\
     a_{21}& a_{22} & \dots & a_{2n}\\
     \vdots&  \vdots &    &  \vdots\\
     a_{m1}& a_{m2}  &\dots & a_{mn}
\end{array}\right) \left(\begin{array}{c}
   x_1 \\
   x_2 \\
   \vdots\\
   x_n 
\end{array}\right)= \left(\begin{array}{c}
   b_1 \\
   b_2 \\
   \vdots\\
   b_n 
\end{array}\right)\]
Es decir, si $A=(a_{ij})$, $x=(x_j)$ y $b=(b_j)$, entonces el sistema anterior es equivalente a $Ax=b$.

\begin{defin}
A la matriz $(A|b)$ se le conoce como la matriz aumentada del sistema de ecuaciones.
\end{defin}
Esta matriz ser\'a \'util en el estudio del n\'umero de soluciones que tiene el sistema y eventualmente determinar el conjunto soluci\'on.

\subsection{Métodos para resolver sistemas de ecuaciones:}

A partir de la matriz aumentada $(A| b)$ se puede resolver al sistema de ecuaciones lineales $Ax = b$, con $A\in \M_{m\times n}$ por lo menos de dos formas que usan del escalonamiento por filas. También se le podrá resolver usando determinantes cuando el sistema tenga igual número de ecuaciones e incógnitas.

La elección de uno u otro método dependerá de las características de tamaño del
sistema a resolver, de sus coeficientes numéricos, de los recursos con que se cuente y,
obviamente, de lo que se pretenda obtener del sistema. Es la práctica la que nos dará
criterios para hacer una buena elección de método.



\begin{itemize}
\item
Eliminación Gaussiana

\begin{enumerate}
\item
Utilizar operaciones filas para reducir la matriz aumentada del sistema a su forma escalonada por filas, hasta obtener la matriz escalonada por filas $(E| g)$ que se asocia al sistema $Ex=g$, que resulta ser equivalente al sistema inicial, es decir ambos sistemas tienen las mismas soluciones.

\item
Despejar el valor de la última incógnita y después utilizar la sustitución hacia atrás para las demás incógnitas.
\end{enumerate}

\item
Eliminación de Gauss Jordan

\begin{enumerate}
\item
Utilizar operaciones filas para reducir la matriz aumentada del sistema a su forma escalonada reducida por filas $(R| s)$.

\item
Si el sistema es consistente, resolver para las variables principales en términos de las variables libres restantes.
\end{enumerate}
\end{itemize}

\begin{rmk}
\begin{enumerate}
    \item Si en el sistema $A·x = b$, con $A\in\M_{m\times n}$, ocurre que su equivalente $(E| g)$ es de rango
$r$ ($r$ filas no nulas), \textbf{habrá $r$ incógnitas dependientes en el sistema y, por ende, $n - r$ incógnitas independientes de las que dependerá el conjunto solución}.
Si resulta que el rango de $(E| g) = n$, que corresponde al n\'umero de variables, entonces todas las incógnitas son
dependientes o determinadas y la solución del sistema $Ax = b$ es única para cada
variable: $x_1= s_1, x_2= s_2, \dots , x_n = s_n$.
\item Si al escalonar $(A| b)$, alguna fila de $(E| g)$ o de $(R| s)$ resulta de la
forma $(0\; 0\; \dots\; 0| g_i \neq 0)$, o $(0\; 0 \;\dots\; 0| s_i \neq 0)$, significa que alguna ecuación del sistema original era contradictoria con las demás puesto que habría una ecuación del tipo $0·x_1 + 0·x_2 + \dots 0·x_n \neq 0$ lo que es imposible. En estos casos, diremos que la solución no existe y que el sistema es
inconsistente o que las ecuaciones son incompatibles. Dicho de otro modo,
la solución será vacía cuando el rango de $A$ sea menor que el rango de $(A|
b)$.
\item Cuando el rango $A$ es igual con el rango de $(A| b)$, podemos concluir que este rango nos indica el número de ecuaciones independientes existentes entre $m$ ecuaciones originales. Todas las filas nulas en $(E| g)$ o $(R| s)$ nos dicen que las ecuaciones asociadas a esas filas eran ecuaciones dependientes o Redundantes, ya que imponían a las variables una condición que ya era exigida por las otras ecuaciones.
\end{enumerate}
Todo lo anterior se puede resumir en el siguiente teorema.
\begin{thm}[Rouché–Frobenius]
Existen soluciones para el sistema si y solo si el rango de la matriz aumentada es igual al rango de la matriz $A$.
Luego,
\[rank(A) < rank(A|b)\Longrightarrow S_S = \emptyset\]
Entonces, si existen soluciones, es decir $rank(A) = rank(A|b)$, tenemos dos situaciones:
\begin{enumerate}
    \item Si el rango de $A$ es $n$, entonces la solución es única. M\'as a\'un, se tiene que $A$ es invertible y $x = A^{-1}b$.
    \item Si $rank(A) = rank(A|b)<n$ existen infinitas soluciones.
\end{enumerate}
\end{thm}


\end{rmk}


\begin{ejer}
Resuelva

\begin{enumerate}
\item
\begin{equation*}
\systeme{
w - x - y + 2z = 1,
2w - 2x - y + 3z = 3,
-w + x - y = -3
}
\end{equation*}


\item
\begin{equation*}
\systeme{
 x_1 - 3x_2 - 5x_3 + x_4 = 4,
 2x_1 + 5x_2 - 2x_3 + 4x_4 = 6
}
\end{equation*}
\end{enumerate}
\end{ejer}
Note que el número de
variables libres será la diferencia entre el total de variables y el rango
de $A$.

\begin{ejer}
Estudiemos los valores del parámetro $\alpha\in \R$ para que la familia de sistemas de ecuaciones lineales:
\begin{equation*}
\systeme{
\alpha x + 3y = -\alpha,
2\alpha x + 5y -\alpha^2 z= 0,
2\alpha x + 5y-z = \alpha +1
}
\end{equation*}
de manera que tenga tenga una solución : i) única ; ii) vacía ; ó, iii) infinitas soluciones.
\end{ejer}




\subsection{Regla de Cramer}

Un sistema de $n$ ecuaciones con $n$ incógnitas: 

$$\begin{matrix} a_{11}x_1 & + & a_{12}x_2 & + & \cdots & + & a_{1n}x_n & = & b_1 \\ a_{21}x_1 & + & a_{22}x_2 & + & \cdots & + & a_{2n}x_n & = & b_2 \\ \vdots & & \vdots & & & & \vdots & & \vdots \\ a_{n1}x_1 & + & a_{n2}x_2 & + & \cdots & + & a_{1n}x_n & = & b_n \end{matrix}$$

se puede escribir de la forma $Ax = b$.


Si $\det A \not= 0$, el sistema tiene solución única. 

Sea $D = \det A$, definamos:


$A_1 = \begin{pmatrix} b_1 & a_{12} & \cdots & a_{1n} \\ b_2 & a_{22} & \cdots & a_{2n} \\ \vdots & \vdots & & \vdots \\ b_n & a_{n2} & \cdots & a_{nn} \end{pmatrix}, A_2 = \begin{pmatrix} a_{11} & b_1 & \cdots & a_{1n} \\ a_{12} & b_2 & \cdots & a_{2n} \\ \vdots & \vdots & & \vdots \\ a_{n1} & b_n & \cdots & a_{nn} \end{pmatrix}, \dots, A_n = \begin{pmatrix} a_{11} & a_{12} & \cdots & b_1 \\ a_{21} & a_{22} & \cdots & b_2 \\ \vdots & \vdots & & \vdots \\ a_{n1} & a_{n2} & \cdots & b_n \end{pmatrix}$.
\\\

y sean $D_1 = \det A_1$, $D_2 = \det A_2$,$\dots, D_n = \det A_n$.


\begin{thm}
Sea $A$ una matriz de $n \times n$ y suponga que $\det A \not= 0$. Entonces la solución única al sistema $Ax = b$ está dada por:

$$x_1 = \frac{D_1}{D}, x_2 = \frac{D_2}{D}, \cdots, x_i = \frac{D_i}{D}, \cdots, x_n = \frac{D_n}{D}$$
\end{thm}


\begin{ejer}
Resuelva el siguiente sistema de ecuaciones utilizando la Regla de Cramer:

\begin{equation*}
\systeme{
2x_1+4x_2+6x_3=18,
4x_1+5x_2+6x_3=24,
3x_1+x_2-2x_3=4
}
\end{equation*}

\end{ejer}

\subsection{Sistemas Lineales Homogéneos}
\begin{defin}
Un sistema de $m \times n$ ecuaciones lineales se llama homogéneo si todas las constantes $b_1,b_2, \dots , b_m$ son cero. Si una o más de las constantes es diferente de cero, decimos que el sistema lineal es no homogéneo.

$$\begin{matrix} a_{11}x_1 &+ & a_{12}x_2 & + & \cdots & + & a_{1n}x_n & = & 0 \\ a_{21}x_1 & + & a_{22}x_2 & + & \cdots & + & a_{2n}x_n & = & 0 \\ \vdots & & \vdots & & &  & \vdots & & \vdots \\ a_{m1}x_1 & + & a_{m2}x_2 & + & \cdots & + & a_{mn}x_n & = & 0 \end{matrix}$$
\end{defin}

\begin{ejem}
\begin{enumerate}
\item
\begin{equation*}
\systeme{
2x_1 + 4x_2 + 6x_3 = 0,
4x_1 + 5x_2 + 6x_3 = 0,
3x_1 + x_2 - 2x_3 = 0
}
\end{equation*}

\item
\begin{equation*}
\systeme{
x_1 + 2x_2 - x_3 = 0,
3x_1 - 3x_2 + 2x_3 = 0,
-x_1 - 11x_2 + 6x_3 = 0
}
\end{equation*}
\end{enumerate}
\end{ejem}


\begin{thm}
Un sistema homogéneo tiene un número infinito de soluciones si tiene más incógnitas que ecuaciones, es decir, si $m < n$.
\end{thm}

\begin{thm}
Sean $x_1$ y $x_2$ dos soluciones al sistema no homogéneo $Ax=b$. Entonces su diferencia $x_1-x_2$ es una solución al sistema homogéneo asociado $Ax= \textbf{0}$.
\end{thm}

\begin{cor}
Sea $x$ una solución particular al sistema no homogéneo y sea $y$ otra solución. Entonces existe una solución $h$ al sistema homogéneo tal que: $y = x + h$.
\end{cor}


\begin{thm}
\textbf{(Resumen)} Sea $A$ una matriz de $n \times n$. Las siguientes afirmaciones son equivalentes:

\begin{enumerate}
\item
$A$ es invertible

\item
La única solución al sistema homogéneo $Ax = 0$ es la solución trivial.

\item
El sistema $Ax = b$ tiene una solución única para cada vector de dimensión $n$ $b$.

\item
$A$ es equivalente por filas a la matriz identidad de $n \times n$.

\item
La forma escalonada por filas de $A$ tiene $n$ pivotes.

\item
$\det A \not = 0$.
\end{enumerate}
\end{thm}

\begin{thm}
Una condición necesaria y suficiente para que $Ax = 0$, sistema de
ecuaciones lineales homogéneas con $n$ variables incógnitas, tenga al menos
una solución diferente de cero es que $rank(A) < n$.
\end{thm}

\textbf{Consecuencias:}
\begin{enumerate}
    \item Si el sistema homogéneo $Ax = 0$, es tal que $rank (A) = n$, o que $|A| \neq 0$, o que $A$ es
invertible, entonces tiene solución trivial es decir $S_H = \{x=(0, \dots, 0) \}$.
\item Si $Ax = 0$ es tal que tiene menos ecuaciones que incógnitas entonces tiene
solucion no trivial, es decir $S_H$ tiene infinitas soluciones.
\item Si $Ax = 0$ es tal que tiene igual número de ecuaciones que de incógnitas entonces:
\begin{enumerate}
    \item Si $|A| \neq 0$ , entonces $S_H = \{0 \}$
    \item  Si $|A| = 0 $, entonces $S_H$ tiene infinitas soluciones.
\end{enumerate}
\item Si $Ax = 0$ tiene al menos una solución no trivial entonces:
\begin{enumerate}
    \item Si $x_0\in S_H$, es decir, es solución del sistema homogéneo, entonces $\lambda x_0 \in S_H$ para todo $\lambda \in K$, pues $A(\lambda x_0) = \lambda(A x_0) = \lambda 0 = 0$.
    \item  Si $x_1 , x_2 \in S_H$ entonces $x_1 + x_2 \in S_H $, pues $Ax_1 = 0$ y $Ax_2 = 0$ entonces
$A(x_1 + x_2 ) = Ax_1 + Ax_2 = 0 + 0 = 0$.
\end{enumerate}
De lo anterior, se concluye que si $Ax = 0$ tiene una solución no trivial entonces
tiene infinitas soluciones.
\end{enumerate}

\section{Diagonalizaci\'on}
Una matriz cuadrada \( A \in \mathbb{R}^{n \times n} \) se dice que es \textbf{diagonalizable} si existe una matriz invertible \( P \in \mathbb{R}^{n \times n} \) y una matriz diagonal \( D \in \mathbb{R}^{n \times n} \) tal que:
\[
A = P D P^{-1}
\]

En los próximos capítulos veremos las condiciones necesarias y suficientes de manera que la matriz sea diagonalizable. Por el momento, sólo describiremos el procedimiento para encontrar las matrices $P$ y $D$.
% Una matriz \( A \) es diagonalizable si y sólo si tiene un conjunto completo de \( n \) vectores propios linealmente independientes, es decir, si existe una base del espacio \( \mathbb{R}^n \) formada por vectores propios de \( A \).
\begin{defin}
    Sea \( A \in \mathbb{R}^{n \times n} \) una matriz cuadrada. Un número \( \lambda \in \mathbb{R} \) se llama \textbf{valor propio} (o \textit{autovalor}) de la matriz \( A \) si existe un vector no nulo \( \vec{v} \in \mathbb{R}^n \), llamado \textbf{vector propio} (o \textit{autovector}), tal que:

\[
A\vec{v} = \lambda \vec{v}
\]

Esta ecuación puede reescribirse como un sistema homogéneo:

\[
(A - \lambda I)\vec{v} = \vec{0}
\]

donde \( I \) es la matriz identidad de orden \( n \). 

Entonces, \( \lambda \) es un valor propio de \( A \) si y sólo si el sistema homogéneo:

\[
(A - \lambda I)\vec{v} = \vec{0}
\]

tiene soluciones no triviales (es decir, \( \vec{v} \neq \vec{0} \)).

Esto ocurre si y sólo si la matriz \( (A - \lambda I) \) es \textbf{invertible}, es decir:

\[
\det(A - \lambda I) = 0
\]

Una vez encontrado un valor propio \( \lambda \), los vectores propios asociados se obtienen resolviendo el sistema lineal homogéneo:

\[
(A - \lambda I)\vec{v} = \vec{0}
\]

\textbf{Nota:} El conjunto de soluciones no triviales de este sistema forma un subespacio vectorial llamado \textit{espacio propio} asociado al valor propio \( \lambda \), pero esto se desarrollará en los capítulos siguientes.
\end{defin}



\subsection*{Procedimiento para Diagonalizar una Matriz}
Dado \( A \in \mathbb{R}^{n \times n} \):

\begin{enumerate}
    \item Calcular los valores propios (\( \lambda \)) resolviendo el \textbf{polinomio característico} dado por:
    \[
    \det(A - \lambda I) = 0
    \]
    
    \item Para cada valor propio \( \lambda_i \), encontrar los vectores propios resolviendo:
    \[
    (A - \lambda_i I)\vec{v} = 0
    \]
    
    \item Si se obtienen \( n \) vectores propios distintos, entonces:
    \[
    P = [\vec{v}_1 \; \vec{v}_2 \; \cdots \; \vec{v}_n], \quad D = \text{diag}(\lambda_1, \lambda_2, \ldots, \lambda_n)
    \]
    
    \item Finalmente:
    \[
    A = P D P^{-1}
    \]
\end{enumerate}

\begin{ejem}
    Sea 
\[
A = \begin{pmatrix}
4 & 1 \\
0 & 2
\end{pmatrix}
\]

\textbf{Paso 1:} Polinomio característico:
\[
\det(A - \lambda I) = \det \begin{pmatrix}
4 - \lambda & 1 \\
0 & 2 - \lambda
\end{pmatrix} = (4 - \lambda)(2 - \lambda)
\]
Valores propios: \( \lambda_1 = 4, \; \lambda_2 = 2 \)

\textbf{Paso 2:} Vectores propios:

Para \( \lambda = 4 \):
\[
(A - 4I)\vec{v} = 0 \Rightarrow \begin{pmatrix}
0 & 1 \\
0 & -2
\end{pmatrix} \vec{v} = 0 \Rightarrow \vec{v}_1 = \begin{pmatrix}1 \\ 0\end{pmatrix}
\]

Para \( \lambda = 2 \):
\[
(A - 2I)\vec{v} = 0 \Rightarrow \begin{pmatrix}
2 & 1 \\
0 & 0
\end{pmatrix} \vec{v} = 0 \Rightarrow \vec{v}_2 = \begin{pmatrix}-1 \\ 2\end{pmatrix}
\]

\textbf{Paso 3:} Matrices \( P \) y \( D \):

\[
P = \begin{pmatrix}
1 & -1 \\
0 & 2
\end{pmatrix}, \quad D = \begin{pmatrix}
4 & 0 \\
0 & 2
\end{pmatrix}
\]

\textbf{Verificación:} \( A = P D P^{-1} \)
\end{ejem}
